\documentclass[border=0pt]{standalone}
% The shared preamble of every docs diagram (docs/diagrams/<name>.tex).
% dvisvgm's driver first, so patterns and opacity survive into the SVG.
\def\pgfsysdriver{pgfsys-dvisvgm.def}
\usepackage[siunitx,RPvoltages]{circuitikz}
\usepackage{xcolor}
% Latin Modern: outline fonts for every glyph a label uses, µ included.
\usepackage[T1]{fontenc}
\usepackage{lmodern}
\usetikzlibrary{arrows.meta,backgrounds,calc,patterns,positioning}
\tikzset{
  % A white card, so diagrams read in mdBook's dark themes too.
  every picture/.append style={show background rectangle,
    background rectangle/.style={fill=white,inner sep=4pt}},
  % Layout geometry: layers told apart by fill and hatching, never colour alone.
  layer a/.style={draw=black,line width=0.6pt,fill=gray!20},
  layer b/.style={draw=black,line width=0.6pt,fill=gray!45},
  layer c/.style={draw=black,line width=0.6pt,fill=white,postaction={pattern=north east lines}},
  layer d/.style={draw=black,line width=0.6pt,fill=white,postaction={pattern=crosshatch dots}},
  marker/.style={draw=black,dashed,line width=0.6pt},
  % A measured distance and its label.
  dim/.style={{Latex[length=2mm]}-{Latex[length=2mm]},line width=0.5pt},
  dim label/.style={font=\small,fill=white,inner sep=1pt},
  % Where a rule fails: thick, dashed and labelled, not only red.
  violation/.style={draw=red!75!black,line width=1.4pt,dash pattern=on 3pt off 2pt},
  violation label/.style={font=\small\bfseries,text=red!75!black},
  caption/.style={font=\small},
}

\begin{document}
\begin{circuitikz}[x=0.15cm,y=0.45cm]
  % Cross-section of the example: n-well 0-60 µm, n-tap 1-3 µm, pdiff 5-55 µm.
  \fill[layer a] (-3,-10) rectangle (63,0);
  \fill[layer b] (0,-5) rectangle (60,0);
  \fill[layer d] (1,-1) rectangle (3,0);
  \fill[layer c] (5,-1) rectangle (55,0);
  \node[caption,anchor=south] at (2,0.2) {ntap};
  \node[caption,fill=white,inner sep=1pt] at (30,-0.5) {pdiff (P+)};
  \node[caption,anchor=west] at (1,-4.2) {nwell};
  \node[caption,anchor=west] at (-2,-9.2) {p substrate};
  % The parasitic PNP under the far end of the diffusion and the well
  % resistance to the tap.
  \node[pnp,anchor=E] (Q) at (50,-1) {};
  \draw (Q.E) -- (50,-1);
  \draw (Q.C) -- (Q.C |- 0,-8.5) node[ground] {};
  \draw (Q.B) to[R, l_={$R_\text{well}$}] (8,0 |- Q.B) -| (2,-1);
  \node[caption,anchor=west] at ($(Q.C)+(1,1)$) {parasitic PNP};
  % The measured distance and the failing point.
  \draw[dim] (3,2.2) -- node[dim label,above] {52 \si{\micro\metre} from the tap, limit 20 \si{\micro\metre}} (55,2.2);
  \draw[densely dotted] (3,0) -- (3,2.2) (55,0) -- (55,2.2);
  \draw[violation] (55,0) circle[radius=6pt];
  \node[violation label,anchor=west] at (57,1) {furthest point};
\end{circuitikz}
\end{document}
